% ==============================================================
% Statistical Inference
% ==============================================================

\section{Statistical Inference}

\subsection{Point estimation: MLE, MOM}

\subsubsection*{Estimators and Estimates}

\begin{definition}
An \textbf{estimator} is a function of sample data used to infer an unknown population parameter $\theta$. Because it depends on random sample data, an estimator $\hat{\theta}$ is itself a \textbf{random variable} with its own probability distribution. The specific numerical value produced from a given sample is called an \textbf{estimate}.
\end{definition}

\begin{example}
The sample mean $\bar{X} = \frac{1}{n}\sum_{i=1}^n X_i$ is an \emph{estimator} of the population mean $\mu$. If a particular sample yields $\bar{x} = 172\,\text{cm}$, then $172\,\text{cm}$ is the \emph{estimate}.
\end{example}

\begin{remark}
Point estimators provide a single value; interval estimators (confidence intervals) provide a range of plausible values. See the next subsection for confidence intervals.
\end{remark}

\subsubsection*{Properties of Good Estimators}

\begin{definition}
Let $\hat{\theta}$ be an estimator of parameter $\theta$, based on a sample of size $n$.
\begin{itemize}
    \item \textbf{Unbiasedness:} $\hat{\theta}$ is \emph{unbiased} if $\mathbb{E}[\hat{\theta}] = \theta$. The \emph{bias} is $\mathrm{Bias}(\hat{\theta}) = \mathbb{E}[\hat{\theta}] - \theta$.
    \item \textbf{Consistency:} $\hat{\theta}$ is \emph{consistent} if $\hat{\theta} \xrightarrow{p} \theta$ as $n \to \infty$; i.e.\ the estimator converges in probability to the true value.
    \item \textbf{Efficiency:} Among all unbiased estimators, $\hat{\theta}$ is \emph{efficient} if it achieves the minimum variance. The \textbf{Cramér--Rao lower bound} provides the theoretical minimum variance for any unbiased estimator.
\end{itemize}
\end{definition}

\begin{example}
The sample mean $\bar{X}$ is unbiased for $\mu$. The sample variance $S^2 = \frac{1}{n-1}\sum_{i=1}^n (X_i - \bar{X})^2$ is unbiased for $\sigma^2$ (the $n-1$ denominator corrects for the bias that arises using $n$).
\end{example}

\subsubsection*{Method of Moments (MOM)}

\begin{definition}
The \textbf{Method of Moments} (MOM) estimates parameters by equating \emph{population moments} to \emph{sample moments} and solving. The $k$-th population moment is $\mu_k' = \mathbb{E}[X^k]$; the corresponding sample moment is $m_k' = \frac{1}{n}\sum_{i=1}^n x_i^k$.
\end{definition}

\begin{example}
For a distribution with mean $\mu$ and variance $\sigma^2$, MOM sets:
\[
\hat{\mu} = m_1' = \bar{x}, \qquad \hat{\sigma}^2 = m_2' - (m_1')^2 = \frac{1}{n}\sum x_i^2 - \bar{x}^2
\]
MOM estimators are simple to compute but may not always be the most efficient.
\end{example}

\subsubsection*{Maximum Likelihood Estimation (MLE)}

\begin{definition}
Given a sample $x_1, \ldots, x_n$ and a parametric model with parameter $\theta$, the \textbf{likelihood function} is:
\[
L(\theta) = \prod_{i=1}^n f(x_i \mid \theta)
\]
The \textbf{maximum likelihood estimator} (MLE) $\hat{\theta}_{\mathrm{MLE}}$ is the value of $\theta$ that maximises $L(\theta)$, or equivalently the \textbf{log-likelihood} $\ell(\theta) = \ln L(\theta)$ (since $\ln$ is monotone):
\[
\hat{\theta}_{\mathrm{MLE}} = \arg\max_{\theta}\; \ell(\theta)
\]
In practice, solve $\frac{\partial \ell}{\partial \theta} = 0$ (the \emph{score equation}).
\end{definition}

\begin{example}
For $X_1,\ldots,X_n \overset{\text{iid}}{\sim} \mathcal{N}(\mu, \sigma^2)$, the log-likelihood is:
\[
\ell(\mu, \sigma^2) = -\frac{n}{2}\ln(2\pi\sigma^2) - \frac{1}{2\sigma^2}\sum_{i=1}^n (x_i - \mu)^2
\]
Setting $\partial\ell/\partial\mu = 0$ gives $\hat{\mu}_{\mathrm{MLE}} = \bar{x}$. Setting $\partial\ell/\partial\sigma^2 = 0$ gives $\hat{\sigma}^2_{\mathrm{MLE}} = \frac{1}{n}\sum(x_i-\bar{x})^2$ (note: biased, unlike MOM's sample variance).
\end{example}

\begin{remark}
MLEs have strong asymptotic properties: they are consistent, asymptotically normal, and asymptotically efficient (achieving the Cramér--Rao bound as $n \to \infty$).
\end{remark}

\subsection{Confidence intervals}

% TODO

\subsection{Hypothesis testing: z, t, chi-squared, F tests}

\subsubsection*{The Chi-Squared ($\chi^2$) Test}

\begin{definition}
The \textbf{chi-squared test statistic} measures how far observed categorical frequencies deviate from expected frequencies:
\[
\chi^2 = \sum_{i} \frac{(O_i - E_i)^2}{E_i}
\]
where $O_i$ is the \emph{observed} count and $E_i$ is the \emph{expected} count in category $i$. Under the null hypothesis, this statistic follows a $\chi^2$ distribution with an appropriate number of degrees of freedom.
\end{definition}

\begin{remark}
A smaller $\chi^2$ value means observed data align closely with expectations (supporting $H_0$). A large $\chi^2$ value — falling in the critical region — provides evidence against $H_0$.
\end{remark}

\subsubsection*{Type I: Goodness-of-Fit Test}

\begin{definition}
A \textbf{goodness-of-fit test} asks: does a sample distribution match a hypothesised population distribution? With $k$ categories, the degrees of freedom are $\nu = k - 1$ (minus one additional for each estimated parameter).
\end{definition}

\begin{example}
A die is rolled 60 times. Under $H_0$ (fair die), each face is expected $E_i = 10$ times. The observed counts are $\{8, 12, 9, 11, 7, 13\}$. Then:
\[
\chi^2 = \frac{(8-10)^2}{10} + \frac{(12-10)^2}{10} + \cdots + \frac{(13-10)^2}{10} = \frac{4+4+1+1+9+9}{10} = 2.8
\]
With $\nu = 5$ degrees of freedom, compare to the critical value $\chi^2_{0.05,5} = 11.07$. Since $2.8 < 11.07$, we fail to reject $H_0$.
\end{example}

\subsubsection*{Type II: Test of Independence}

\begin{definition}
A \textbf{test of independence} determines whether two categorical variables are associated. Data are arranged in an $r \times c$ contingency table. The expected count for cell $(i,j)$ is:
\[
E_{ij} = \frac{(\text{row } i \text{ total}) \times (\text{column } j \text{ total})}{\text{grand total}}
\]
Degrees of freedom: $\nu = (r-1)(c-1)$.
\end{definition}

\begin{example}
To test whether gender and voting preference are independent, arrange counts into a $2 \times 3$ contingency table (2 genders, 3 parties). Compute $E_{ij}$ for each cell, then sum $\frac{(O_{ij}-E_{ij})^2}{E_{ij}}$ across all cells. With $\nu = (2-1)(3-1) = 2$, compare to $\chi^2_{0.05,2} = 5.99$.
\end{example}

\subsubsection*{Requirements and Limitations}

\begin{remark}
The $\chi^2$ test requires:
\begin{itemize}
    \item \textbf{Categorical data} — it does not apply to continuous measurements directly.
    \item \textbf{Independent observations} — each subject appears in exactly one cell.
    \item \textbf{Sufficient expected counts} — every cell should have $E_i \geq 5$. If this fails, merge categories or use Fisher's exact test.
\end{itemize}
The $\chi^2$ distribution is an approximation; it becomes accurate as sample size grows. For a $2 \times 2$ table with small samples, apply Yates' continuity correction:
\[
\chi^2 = \sum_{i} \frac{(|O_i - E_i| - 0.5)^2}{E_i}
\]
\end{remark}

\begin{remark}
For an application of the goodness-of-fit test to real-world fraud detection, see the Benford's Law subsection below, where the $\chi^2$ statistic is used to detect suspiciously uniform leading-digit distributions in financial data.
\end{remark}

\subsubsection*{The $F$-Test}

\begin{definition}[$F$-test]
An \textbf{$F$-test} is used to test joint hypotheses about multiple parameters simultaneously. Under the null, the test statistic follows an $F(q, n-k)$ distribution where $q$ is the number of restrictions and $n-k$ is the residual degrees of freedom.
\end{definition}

\begin{remark}[Why we need $F$ instead of repeated $t$-tests]
Testing $H_0\colon \beta_1 = 0$ \emph{and} $\beta_2 = 0$ with two separate $t$-tests overstates confidence: each test has a false-positive probability, and running two inflates the overall Type I error rate. The $F$-test tests both restrictions simultaneously and controls the error rate correctly.
\end{remark}

\begin{theorem}[Large-sample $F$-statistic for $q$ restrictions]
In large samples, when $q$ restrictions $H_0\colon R\beta = r$ are tested jointly, the $F$-statistic follows
\[
F \sim F(q, \infty).
\]
This is because each restricted $t$-statistic is approximately $N(0,1)$ under $H_0$, so their squares are $\chi^2(1)$ and their sum is $\chi^2(q)$:
\[
F = \frac{\chi^2(q)}{q} \sim F(q,\infty).
\]
\end{theorem}

\begin{remark}[Connection to the distributions in ch.\ 62]
$F(q,\infty)$ is the limit of $F(q, d_2)$ as $d_2 \to \infty$, at which point the denominator stabilises to 1 and $F = \chi^2(q)/q$. For regression in finite samples, the denominator uses the residual $\chi^2(n-k)$ in place of the infinite-df limit.
\end{remark}

\subsection{Benford's Law and Fraud Detection}

\begin{definition}
\textbf{Benford's Law} (also called the \textit{First-Digit Law}) states that in many naturally occurring datasets, the probability that the leading (first significant) digit is $d \in \{1, 2, \ldots, 9\}$ follows a logarithmic distribution:
\[
P(D = d) = \log_{10}\!\left(1 + \frac{1}{d}\right)
\]
\end{definition}

\begin{remark}
The predicted probabilities are approximately:
\begin{center}
\begin{tabular}{c|ccccccccc}
Digit $d$ & 1 & 2 & 3 & 4 & 5 & 6 & 7 & 8 & 9 \\
\hline
$P(D=d)$ & 30.1\% & 17.6\% & 12.5\% & 9.7\% & 7.9\% & 6.7\% & 5.8\% & 5.1\% & 4.6\%
\end{tabular}
\end{center}
Crucially, the digit \textbf{1} appears as the leading digit nearly a third of the time --- far more often than a naive uniform assumption ($11.\overline{1}\%$) would suggest.
\end{remark}

\subsubsection*{Logarithmic Interpretation}

The law arises from the \emph{gaps} between consecutive base-10 logarithms of integers:
\[
P(D = d) = \log_{10}(d+1) - \log_{10}(d) = \log_{10}\!\left(\frac{d+1}{d}\right) = \log_{10}\!\left(1+\frac{1}{d}\right)
\]
These gaps are largest at $d=1$ and shrink monotonically toward $d=9$, explaining why lower digits dominate.

\begin{theorem}[Condition for Benford's Law]
Benford's Law holds for a dataset when the \textbf{fractional part} of $\log_{10}(x)$ --- i.e.\ $\{\log_{10}(x)\}$ --- is uniformly distributed on $[0,1)$. Under this condition, the probability of each leading digit is exactly the logarithmic gap above.
\end{theorem}

\subsubsection*{Why Naturally Occurring Data Follows This Pattern}

Consider a quantity growing multiplicatively (e.g.\ compound interest at 10\% per year):
\begin{itemize}
    \item Going from \$1 to \$2 requires a $100\%$ increase (and takes many years).
    \item Going from \$8 to \$9 requires only an $11\%$ increase (reached quickly).
\end{itemize}
Numbers therefore spend \emph{proportionally more time} with leading digit 1 than with 9, regardless of the unit of measurement.

\begin{theorem}[Scale Invariance]
Benford's Law is \textbf{scale invariant}: multiplying an entire dataset by any positive constant (e.g.\ converting currencies or units) does not change the leading-digit distribution. This is a hallmark property that distinguishes naturally generated data from fabricated data.
\end{theorem}

\subsubsection*{Limitations}

Benford's Law \emph{fails} when:
\begin{itemize}
    \item The data spans a narrow range of magnitudes (e.g.\ adult human heights in metres, all clustered near $1$--$2$ m).
    \item The dataset is artificially constrained (e.g.\ prices set to always begin with \$9.99).
    \item Data are purely random (uniformly or normally distributed with no multiplicative structure).
\end{itemize}

\subsubsection*{Application: Fraud Detection}

\begin{example}
Given an unlabelled dataset of financial figures, one can test conformity to Benford's Law using a \textbf{chi-squared goodness-of-fit test}:
\[
\chi^2 = \sum_{d=1}^{9} \frac{(O_d - E_d)^2}{E_d}, \qquad E_d = N \cdot \log_{10}\!\left(1+\frac{1}{d}\right)
\]
where $O_d$ is the observed count of leading digit $d$ and $N$ is the total number of values. A statistically significant deviation (especially a distribution that is \emph{too uniform}) is a red flag for fabricated data, since humans inventing numbers tend to distribute digits more evenly than nature does.
\end{example}

\begin{remark}
Benford's Law has been used to audit elections, detect accounting fraud, and expose billion-dollar financial scandals. It is context-free: no knowledge of what the numbers represent is required, only that the data plausibly spans several orders of magnitude. Non-conformity is not definitive proof of fraud, but it is a statistically grounded signal warranting further investigation.
\end{remark}
